\scanpage{163}

On dit qu'un graphe à homothéties est \uline{régulier} si $G_0$ (le groupe de ses autom \uline{stricts}) est transitif sur l'ens $\overrightarrow{A_0}$ de ses arêtes orientées, \inserted{et si $(S,A_0)$ n'a pas de sommet isolé et est connexe.} On dit qu'un multigraphe \deleted{est} à homothéties est régulier, si les graphes constituants le sont.

Si \deleted{$\unsure{G_0|\mathcal G}$} $(S,A_0,Z)$ est \inserted{un gr. à hom.} \uline{régulier}, $G_0$ étant transitif sur $S$, $Z$ opère \uline{librement} sur $S$. En tous cas, j'appelle \uline{\unsure{vagons|wagons|voyage}} du graphe à homothéties, \deleted{les} orbites \deleted{pour} \inserted{de} $Z$ dans $S$.

Si \inserted{on a un multigr. à hom.} \deleted{\unsure{$($}} $(S,\pi,Z)$ \deleted{est} régulier, alors \deleted{pour $i,j\in\unsure{R_\pi|\pi|\Pi}$ distincts, on a}

\[ A_i\cap A_j=\varnothing\qquad\text{si }i,j\in\pi,\quad i=j \tag*{(14)} \]

car si on a $a\in A_i\deleted{\cap A_j}$, alors $A_i$ est déterminé par $a$ comme l'orbite de $a\in\unsure{\mathcal P|\mathfrak P}_2(S)$ sous $G_0$. On trouve donc une application surjective

% slanted
\margin{La structure \inserted{de multigraphe rég.} $(S,\pi,Z)$ équivaut alors à la structure $(S,\unsure{A|F},R,Z)$, où $A\subset\mathcal P_2(S)$ muni de $Z\subset\unsure{\mathfrak S|\mathscr S}_S$, une rel. d'\unsure{équiv.|équiva.} $R$, \unsure{et où|et} (plus \unsure{ou moins|comme} bien sûr, savoir que chacun des $(S,A_i,Z)$ forme un ``graphe'' \deleted{\unsure{une}} str. de graphes à homothéties \uline{régulières}\ldots}

\[ A\overset{\mathrm{déf}}{=}\bigcup_{i\in\pi} A_i\xrightarrow{\tau_A}\pi \tag*{(15)} \]

associée à une partition

\[ A=\coprod_{i\in\pi} A_i. \tag*{(16)} \]

\deleted{Soit. On dit qu'un} \unsure{Revenant|Revenons} à un $A_0$, et posons

\[
\begin{aligned}
 F_0&=\{f\in\mathcal P_3(S)\mid\forall a\in\mathcal P_2(f),\ a\in A_0\}\\
 \text{d'où }\quad A_0&=\{a\in\mathcal P_2(S)\mid\exists f\in F_0,\ \text{\unsure{t.|,}}~ f\supset a\}.
\end{aligned}\tag*{(17)}
\]

Une structure de graphe sera dite \deleted{\unsure{prétri}} \deleted{``triangulable''} \uline{``prépolyédrale''}, si

\[ P_0=(S,A_0,F_0,\text{relations d'incidence \unsure{évidentes|évidente}}) \tag*{(18)} \]

% archive sheet number: 162
